\chapter{序論}

\section{研究背景}

近年、フィールドロボットは、農業、建設、災害対応、環境保全、惑星探査など、様々な分野において重要な役割を果たしている。これらのロボットは、人間にとって危険または困難な作業を自動化し、効率性と安全性を大幅に向上させる。特に、不整地環境における自律移動技術は、フィールドロボットの実用化において最も重要な要素の一つとなっている。不整地環境とは、海岸、山地、森林、災害現場、農地など、地形が複雑で予測困難な環境を指す。

不整地環境では、急斜面、凹凸、表面粗さ、障害物、可変的な路面状態などが複雑に混在している。これらの地形要因は、ロボットの走行安定性、消費エネルギー、移動速度に大きな影響を与える。例えば、急斜面を直線的に登る経路は最短距離であるが、ロボットが転倒や滑落するリスクが高く、また、モーターやバッテリーに大きな負荷をかける。一方、緩やかな斜面を迂回する経路は距離が長くなるが、安全性が高く、エネルギー効率も良好である。このように、不整地環境における経路計画では、単に距離を最小化するだけでなく、地形の特性を考慮した総合的な評価が必要となる。

従来の経路計画手法（A*、RRT*、D*など）は、主に経路長や計算効率に焦点を当てており、グリッドベースの環境表現において高い性能を達成する。これらの手法は、屋内環境や平坦な環境では有効であるが、不整地環境における地形の複雑さを十分に考慮していない。その結果、アルゴリズムが生成する経路は、数学的には最適であっても、実際のロボット走行においては、急斜面や粗い地形を通過するため、走行困難または危険となる場合がある。特に、地形情報をコスト関数に明示的に組み込まない場合、経路計画アルゴリズムは、地形の良し悪しを判断できず、最短経路が常に最良の経路として選択されてしまう。

さらに、フィールドロボットは、限られた計算資源と通信帯域の下で動作することが多く、リアルタイム性が求められる。複雑な環境認識や経路計画アルゴリズムは、高い計算コストを要するため、実機での実行可能性を考慮したアルゴリズム設計が不可欠である。また、環境が動的に変化する場合（新たな障害物の出現、天候による地形変化など）には、経路の再計画を高速に行う必要がある。このような背景から、計算効率を維持しながら地形情報を考慮できる経路計画手法の開発が強く求められている。

\section{関連研究}

\subsection{A*アルゴリズムとその派生手法}

A*アルゴリズムは、ヒューリスティック関数を用いて探索を効率化する代表的な経路探索手法である。A*は、グリッドベースの環境において、最短経路を保証しつつ、効率的に探索を行うことができる。評価関数$f(n) = g(n) + h(n)$を用いて、始点からの実コスト$g(n)$とゴールまでの推定コスト$h(n)$を組み合わせることで、最適な経路を発見する。ヒューリスティック関数$h(n)$が許容的（真のコストを過大評価しない）である場合、A*は最適解を保証する。

A*の派生手法として、サンプリングベースのRRT*（Rapidly-exploring Random Tree*）、動的環境に対応するD*Lite、階層的探索を行うHPA*（Hierarchical Pathfinding A*）などが提案されている。RRT*は、ランダムサンプリングにより高次元空間での経路計画が可能であり、複雑な環境での経路探索に適している。D*Liteは、環境の変化に応じて効率的に経路を更新できるため、動的環境での利用に優れる。HPA*は、環境を階層的に分割することで、大規模環境での探索を高速化する。

しかし、これらの手法は、主に経路長や計算効率に焦点を当てており、地形の複雑さや走行困難性を明示的に考慮していない。そのため、不整地環境では、生成された経路が実際のロボット走行に適していない場合がある。本研究では、A*を基盤として、地形情報を統合したTerrain-Aware A*を提案する。

\subsection{Field D*アルゴリズム}

Field D*は、D*アルゴリズムの拡張版であり、隣接セル間で補間を行うことで、グリッドベースの制約を緩和する。通常のグリッドベースのアルゴリズムでは、経路がグリッドのエッジに沿って生成されるため、不自然な折れ曲がりが生じる。Field D*は、セル内の任意の点を経由することで、より直線的で滑らかな経路を生成できる。この補間技術により、経路長の短縮とロボットの動作の滑らかさが向上する。

Field D*は、動的環境における再計画能力にも優れており、環境の変化に応じて経路を効率的に更新できる。これは、D*が持つインクリメンタルな探索手法を継承しているためである。Field D*は、火星探査ローバーMars Exploration Roverなどの実際のフィールドロボットにも適用されており、その有効性が実証されている。特に、岩石の多い火星表面において、Field D*は滑らかで効率的な経路を生成し、ローバーの移動距離を短縮することに成功した。

本研究では、Field D*の補間技術を基盤として、地形情報を統合したField D* Hybridを開発する。Field D* Hybridは、補間による経路短縮と、地形考慮による安全性向上を同時に実現することを目指す。

\subsection{地形考慮型経路計画}

地形を考慮した経路計画に関する研究も存在する。これらの研究では、地形の傾斜、障害物密度、路面状態などをコスト関数に統合し、走行しやすい経路を生成する。例えば、Mars探査ローバーの経路計画では、岩石の密度や斜面の傾斜を考慮した手法が用いられている。ローバーは、急斜面や岩石の多い領域を避けることで、安全に目的地に到達できる。

また、農業ロボットでは、畝や溝を避けるために、地形情報を利用した経路計画が研究されている。農業ロボットは、作物を傷つけずに移動する必要があるため、畝を地形の悪い領域として扱い、これを回避する経路を生成する。さらに、軍事用の無人地上車両（UGV）では、敵の視界を避けるために、地形の起伏を利用した経路計画が研究されている。

しかし、多くの先行研究は、特定の地形タイプ（例：傾斜のみ、障害物密度のみ）に限定されており、複数の地形要因（標高差、傾斜、表面粗さ）を統合的に評価する手法は少ない。また、RTAB-Mapのようなリアルタイム点群取得システムと統合し、実装可能性の高いシステムとして提案されている研究も限定的である。本研究では、複数の地形要因を統合した地形複雑度を定義し、RTAB-Mapというリアルタイム3次元SLAMツールと統合することで、実装可能性の高い地形考慮型経路計画システムを提案する。

\subsection{RTAB-Mapによる環境認識}

環境認識技術として、RTAB-Map（Real-Time Appearance-Based Mapping）が広く用いられている。RTAB-Mapは、RGB-DカメラやLiDARから得られるセンサーデータを用いて、リアルタイムで3次元環境地図を生成するSLAMライブラリである。RTAB-Mapは、視覚的な特徴点とループクロージャ検出を組み合わせることで、大規模環境においても長時間にわたり安定して動作する。

RTAB-Mapは、メモリ管理機構を持ち、長時間の運用でもメモリ使用量を一定に保つことができる。これは、古いデータを選択的に削除し、重要なデータのみを保持することで実現される。また、RTAB-Mapは、ループクロージャ検出により、累積誤差を修正し、高精度な地図を生成できる。これらの特徴により、RTAB-Mapは、フィールドロボットの環境認識に広く利用されている。

本研究では、RTAB-Mapを用いて取得する点群データから地形情報を抽出し、経路計画に活用する手法を提案する。点群データは、環境の3次元構造を詳細に表現しており、標高差、傾斜、表面粗さなどの地形特性を定量的に評価することができる。

\section{研究目的}

不整地環境における自律移動ロボットの実用化には、以下の3つの課題が存在する：(1) 地形の複雑さを定量評価する統一指標の不在、(2) 経路長と安全性のトレードオフの定量的理解の不足、(3) 実時間性と信頼性を両立する手法の不足。従来の経路計画手法は、経路長の最適化に焦点を当てており、地形特性（傾斜、粗さ、標高差）を明示的に考慮していない。そのため、生成される経路は数学的には最適でも、実際のロボット走行においては危険または非効率となる場合がある。

本研究は、これらの課題を解決するため、地形考慮型経路計画フレームワークを構築し、以下を達成する。

\textbf{（1）地形複雑度の定量評価フレームワークの構築：}
RTAB-Mapから得られる点群データから、標高差、傾斜角、表面粗さを統合した地形複雑度指標を構築する。点群データを解像度0.2mのボクセルグリッドに変換し、各ボクセルにおいて地形複雑度を計算する。予備実験により最適な重みパラメータ（$\alpha=\beta=0.4$, $\gamma=0.2$）を決定し、0から1の範囲に正規化された地形コストマップを経路計画アルゴリズムに提供する。これにより、実時間での地形評価を可能にする。

\textbf{（2）地形考慮型経路計画アルゴリズムの開発と実用化：}
Terrain-Aware A*（TA*）により地形回避経路を生成し、96シナリオの大規模ベンチマークにより性能を統計的に検証する。さらに、Field D*の補間技術を統合したField D* Hybridを開発し、実時間対応（計算時間0.175秒、成功率100\%）を実現する。比較実験では、4つの経路計画手法の性能を、経路長、計算時間、成功率、地形コストの観点から定量的に比較する。

\textbf{（3）実用化に向けた性能特性の解明：}
地形複雑度別（低<0.15、中0.15--0.55、高$\geq$0.55）の性能分析により、各手法の適用条件と限界を明確化する。失敗シナリオの分析により、実装上の注意点を特定し、実環境への展開指針を提示する。これにより、農業ロボット、災害対応ロボット、惑星探査ローバーなど、多様な分野への応用可能性を示す。

本研究により、不整地環境におけるロボットの走行安定性と安全性を大幅に向上させる経路計画手法を実現する。地形の悪い領域を回避することで、転倒や滑落のリスクを低減し、消費エネルギーを削減できる。提案するフレームワークは、実用レベルの性能（計算時間0.175秒、成功率100\%）を達成しており、実環境への即座の実装が可能である。

\section{本論文の構成}

本論文の構成は以下の通りである。

第2章では、提案手法の将来的な実装を見据えたハードウェア構成とソフトウェア環境について述べる。本研究ではシミュレーション環境での評価を行ったが、実装可能性を示すため、想定するロボットシステム（AgileX社製クローラー型UGV「BUNKER」とセンサー構成）について説明する。また、将来的に使用予定のROS 2による開発環境についても述べる。

第3章では、環境認識と地形情報取得の手法について述べる。自己位置推定のためのセンサーフュージョン（RTK-GPS、wheelオドメトリ、IMU）、RTAB-Mapによる3次元環境地図生成、点群データからの地形情報抽出手法について詳しく説明する。ボクセルグリッド生成、地形複雑度の計算、地形コストマップの生成について述べる。

第4章では、提案手法であるTerrain-Aware A*（TA*）アルゴリズムとField D* Hybridについて詳しく説明する。まず、A*アルゴリズムの基本原理を述べた後、既存手法の問題点と本研究の独自性を明確化する。次に、TA*の拡張評価関数と地形コストの計算方法を説明し、パラメータ調整の戦略について述べる。最後に、Field D* Hybridとの統合方法とアルゴリズムの全体フローを説明する。

第5章では、実験設定とベンチマークデータセットについて述べる。Dataset3の生成方法、構成（basic、obstacles、terrain）、各カテゴリの特徴を説明する。また、比較に用いる4つの経路計画手法と、評価指標（経路長、計算時間、成功率、地形コスト）について詳しく述べる。実験環境（Ubuntu 22.04、Python 3.10）とパラメータ設定についても説明する。

第6章では、実験結果を報告する。各手法の経路長、計算時間、成功率を比較し、地形評価結果を詳しく分析する。特に、代表的なシナリオにおける提案手法の性能を詳細に評価し、経路長と地形品質のトレードオフを明らかにする。

第7章では、実験結果に対する考察を行う。提案手法の性能評価、2つの評価指標（点平均法、セグメント加重法）の使い分け、実用化に向けた課題について議論する。また、適用シーンや今後の改善方向についても述べる。

第8章では、今後の展望について述べる。アルゴリズムの拡張（動的環境対応、マルチモーダル評価、機械学習の統合）、実環境実験の必要性、他分野への応用可能性（農業ロボット、災害対応ロボット、惑星探査ローバー）について述べる。

第9章では、本研究の結論をまとめる。研究目的の再確認、主要な成果の列挙、今後の展望を簡潔に述べ、本論文を締めくくる。
